
\section*{\underline{Exposé VI}\\[3pt] \underline{\underline{LE \(K^{\bullet}\) D'UN FIBRÉ PROJECTIF :}}\\[3pt] \underline{\underline{CALCULS ET CONSÉQUENCES}}}
% Exposé VI: the print's doubly underlined S and underlined Proj (scan check, 10 Nov)
\renewcommand{\Sym}{\mathbf{S}}
\renewcommand{\Projz}{\underline{\operatorname{Proj}}_0}
\begin{center}
\underline{par Pierre BERTHELOT}
\end{center}

\section*{1. Calcul du \(K^{\bullet}\) d'un fibré projectif : cas des faisceaux localement libres de type fini}

Dans ce paragraphe, la notation \(K^{\bullet}(X)\) désignera le groupe de Grothendieck des faisceaux localement libres de type fini sur un schéma \(X\).

Le but du paragraphe est de démontrer le théorème suivant :

\textbf{Théorème 1.1.} \uline{Soient} \(S\) \uline{un schéma quasi-compact}, \(\underline{E}\) \uline{un} \(\underline{O}_S\)\uline{-Module localement libre de rang} \(r+1\), \(X=\mathbf P(\underline{E})\) \uline{le fibré projectif associé} ; \uline{le morphisme structural} : \(X\xrightarrow{f}S\) \uline{définit un homomorphisme d'anneaux} : \(K^{\bullet}(S) \xrightarrow{f^*} K^{\bullet}(X)\) \uline{munissant} \(K^{\bullet}(X)\) \uline{d'une structure de} \(K^{\bullet}(S)\)-\uline{algèbre}. \uline{Soit} \(\xi\in K^{\bullet}(X)\) \uline{la classe du faisceau} \(\underline{O}_X(1)\). \uline{Alors} \(K^{\bullet}(X)\) \uline{est un} \(K^{\bullet}(S)\)-\uline{module libre admettant pour base} \(1,\ldots,\xi^r\).

\uline{Plan de la démonstration} :

\(1^\circ\)) Les puissances de \(\xi\) engendrent \(K^{\bullet}(X)\) comme \(K^{\bullet}(S)\)-module.

Pour le montrer, nous allons montrer les points suivants :

\begin{enumerate}[label=\alph*)]
\item Si \(\underline{F}\) est un \(\underline{O}_X\)-Module localement libre de type fini, il existe un \(\Sym(\underline{E})\)-Module gradué \(\underline{M}\) de présentation finie, plat sur \(\underline{O}_S\), tel que \(\underline{F}=\Projz(\underline{M})\) (\(\Sym(\underline{E})\) étant l'algèbre symétrique de \(\underline{E}\)~; pour la notation \(\Projz\), voir EGA II 8.12.1).

\item Il existe une résolution graduée finie de \(\underline{M}\) par des \(\Sym(\underline{E})\)-Modules gradués \(\underline{L}_i\) localement libres de type fini : \(0\longrightarrow \underline{L}_{r+1}\longrightarrow\ldots\longrightarrow \underline{L}_0\longrightarrow \underline{M}\longrightarrow 0\).

\item Les \(\underline{L}_i\) admettent une filtration finie dont les quotients sont du type \(\underline{N}_j\otimes_{\underline{O}_S}\Sym(\underline{E})(j)\), les \(\underline{N}_j\) étant localement libres de type fini sur \(S\).
\end{enumerate}

